\section{微分形式}\label{sec:differential-forms}

交错张量场是非常重要的对象, 称为\textbf{形式}.

\begin{RiemannDefinition}\label{def:differential-form}
设 $M$ 是光滑流形. $M$ 上的\textbf{$k$ 次形式} (或 \textbf{$k$-形式}) 是定义在 $M$ 上的交错 $k$-张量场, 即给每个点 $p\in M$ 指派一个元素 $\omega_p\in\Lambda^k(T_p^*M)$ 的映射 $\omega$.
\end{RiemannDefinition}

$M$ 上的 $k$-形式空间显然是一个向量空间. 由 Theorem~\ref{thm:basis-and-dimension-of-alternating-tensors}, 给定坐标系 $x:W\to\mathbb{R}^n$, $W$ 上任意 $k$-形式都可写成
\[
  \omega=\sum_I\omega_I\,\mathrm{d}x^I
\]
其中 $I=(i_1,\ldots,i_k)$ 表示由 $\{1,\ldots,n\}$ 中的整数组成的任意递增指标序列, $\mathrm{d}x^I$ 是形式 $\mathrm{d}x^{i_1}\wedge\cdots\wedge\mathrm{d}x^{i_k}$, 而 $\omega_I$ 是定义在 $W$ 上的函数. 容易验证, $\omega$ 在基 $\{\mathrm{d}x^{i_1}\otimes\cdots\otimes\mathrm{d}x^{i_k}\}$ 下的分量为 $\pm\omega_I$. 因此, 若对极大图册中的所有坐标系, 函数 $\omega_I$ 都是光滑的, 则 $\omega$ 是可微 $(k,0)$-张量 (此时称它为\textbf{微分形式}). $M$ 上的微分 $k$-形式集合记为 $\Omega^k(M)$. 以下用 \textquotedblleft 形式\textquotedblright 一词指微分形式.

给定可微流形之间的光滑映射 $f:M\to N$, 可以用拉回运算从 $N$ 上的形式诱导出 $M$ 上的形式 (参见 Definition~\ref{def:pullback-of-covariant-tensor-field}), 因为交错张量场的拉回仍是交错张量场.

\begin{RiemannRemark}\label{rem:pullback-of-zero-form}
若 $g:N\to\mathbb{R}$ 是 0-形式, 即函数, 则其拉回定义为 $f^*g=g\circ f$.
\end{RiemannRemark}

容易验证, 形式的拉回满足下列性质.

\begin{RiemannProposition}\label{prop:pullback-of-differential-forms}
设 $f:M\to N$ 是可微映射, $\alpha,\beta$ 是 $N$ 上的形式. 则
\begin{enumerate}[label=(\roman*),leftmargin=*]
\item $f^*(\alpha+\beta)=f^*\alpha+f^*\beta$;
\item 对任意函数 $g\in C^\infty(N)$, 有 $f^*(g\alpha)=(g\circ f)f^*\alpha=(f^*g)(f^*\alpha)$;
\item $f^*(\alpha\wedge\beta)=(f^*\alpha)\wedge(f^*\beta)$;
\item 对任意映射 $g\in C^\infty(L,M)$, 有 $g^*(f^*\alpha)=(f\circ g)^*\alpha$, 其中 $L$ 为可微流形.
\end{enumerate}
\end{RiemannProposition}

\begin{RiemannProof}
见 Exercise~\ref{ex:differential-forms}(\ref{ex:prove-pullback-of-differential-forms}).
\end{RiemannProof}

\begin{RiemannExample}\label{ex:coordinate-pullback-of-differential-form}
若 $f:M\to N$ 可微, 并考虑 $M$ 和 $N$ 上各自的坐标系 $x:V\to\mathbb{R}^m$, $y:W\to\mathbb{R}^n$, 则对 $i=1,\ldots,n$, 有 $y^i=\widehat f^i(x^1,\ldots,x^m)$, 其中 $\widehat f=y\circ f\circ x^{-1}$ 是 $f$ 的局部表示. 若 $\omega=\sum_I\omega_I\,\mathrm{d}y^I$ 是 $W$ 上的 $k$-形式, 则由 Proposition~\ref{prop:pullback-of-differential-forms},
\begin{align*}
  f^*\omega
  &=f^*\left(\sum_I\omega_I\,\mathrm{d}y^I\right)
   =\sum_I(f^*\omega_I)(f^*\mathrm{d}y^I)\\
  &=\sum_I(\omega_I\circ f)(f^*\mathrm{d}y^{i_1})\wedge\cdots\wedge(f^*\mathrm{d}y^{i_k}).
\end{align*}
此外, 对 $v\in T_pM$,
\[
  (f^*(\mathrm{d}y^i))_p(v)
  =(\mathrm{d}y^i)_{f(p)}((\mathrm{d}f)_pv)
  =\bigl(\mathrm{d}(y^i\circ f)\bigr)_p(v),
\]
即 $f^*(\mathrm{d}y^i)=\mathrm{d}(y^i\circ f)$. 因此,
\begin{align*}
  f^*\omega&=\sum_I(\omega_I\circ f)\,\mathrm{d}(y^{i_1}\circ f)\wedge\cdots\wedge\mathrm{d}(y^{i_k}\circ f)\\
  &=\sum_I(\omega_I\circ f)\,\mathrm{d}(\widehat f^{i_1}\circ x)\wedge\cdots\wedge\mathrm{d}(\widehat f^{i_k}\circ x).
\end{align*}
若 $k=\dim M=\dim N=n$, 则拉回 $f^*\omega$ 很容易由 Theorem~\ref{thm:pullback-top-degree-determinant} 算出, 由该定理有
\begin{equation}\label{eq:pullback-top-form-in-coordinates}
  (f^*(\mathrm{d}y^1\wedge\cdots\wedge\mathrm{d}y^n))_p
  =\det(\mathrm{d}\widehat f)_{x(p)}(\mathrm{d}x^1\wedge\cdots\wedge\mathrm{d}x^n)_p.
\end{equation}
\end{RiemannExample}

给定 $M$ 上任意形式 $\omega$ 和参数化 $\varphi:U\to M$, 可以取 $\omega$ 沿 $\varphi$ 的拉回, 得到一个定义在开集 $U$ 上的形式, 称为 $\omega$ 在该参数化下的\textbf{局部表示}.

\begin{RiemannExample}\label{ex:local-representation-of-coordinate-one-form}
设 $x:W\to\mathbb{R}^n$ 是光滑流形 $M$ 上的坐标系, 考虑定义在 $W$ 上的 1-形式 $\mathrm{d}x^i$. 沿相应参数化 $\varphi:=x^{-1}$ 的拉回 $\varphi^*\mathrm{d}x^i$ 是 $\mathbb{R}^n$ 的开子集 $U$ 上的 1-形式, 满足
\begin{align*}
  (\varphi^*\mathrm{d}x^i)_x(v)
  &=(\varphi^*\mathrm{d}x^i)_x\left(\sum_{j=1}^nv^j\left(\frac{\partial}{\partial x^j}\right)_x\right)\\
  &=(\mathrm{d}x^i)_p\left(\sum_{j=1}^nv^j(\mathrm{d}\varphi)_x\left(\frac{\partial}{\partial x^j}\right)_x\right)\\
  &=(\mathrm{d}x^i)_p\left(\sum_{j=1}^nv^j\left(\frac{\partial}{\partial x^j}\right)_p\right)
  =v^i=(\mathrm{d}x^i)_x(v),
\end{align*}
其中 $x\in U$, $p=\varphi(x)$, 且 $v=\sum_{j=1}^nv^j\bigl(\frac{\partial}{\partial x^j}\bigr)_x\in T_xU$. 因此, 如同已有 $\bigl(\frac{\partial}{\partial x^i}\bigr)_p=(\mathrm{d}\varphi)_x\bigl(\frac{\partial}{\partial x^i}\bigr)_x$, 现在有 $(\mathrm{d}x^i)_x=\varphi^*(\mathrm{d}x^i)_p$, 所以 $(\mathrm{d}x^i)_p$ 是 $W$ 上以 $(\mathrm{d}x^i)_x$ 为其在 $U$ 上局部表示的 1-形式.
\end{RiemannExample}

若 $\omega=\sum_I\omega_I\,\mathrm{d}x^I$ 是定义在 $\mathbb{R}^n$ 的某个开子集上的 $k$-形式, 则如下定义一个称为 $\omega$ 的\textbf{外微分}的 $(k+1)$-形式:
\[
  \mathrm{d}\omega:=\sum_I\mathrm{d}\omega_I\wedge\mathrm{d}x^I.
\]

\begin{RiemannExample}\label{ex:exterior-derivative-of-angular-form}
考虑定义在 $\mathbb{R}^2\setminus\{0\}$ 上的形式 $\omega=-\frac{y}{x^2+y^2}\,\mathrm{d}x+\frac{x}{x^2+y^2}\,\mathrm{d}y$. 则
\begin{align*}
  \mathrm{d}\omega
  &=\mathrm{d}\left(-\frac{y}{x^2+y^2}\right)\wedge\mathrm{d}x
    +\mathrm{d}\left(\frac{x}{x^2+y^2}\right)\wedge\mathrm{d}y\\
  &=\frac{y^2-x^2}{(x^2+y^2)^2}\,\mathrm{d}y\wedge\mathrm{d}x
    +\frac{y^2-x^2}{(x^2+y^2)^2}\,\mathrm{d}x\wedge\mathrm{d}y=0.
\end{align*}
\end{RiemannExample}

外微分满足下列性质:

\begin{RiemannProposition}\label{prop:properties-of-exterior-derivative}
若 $\alpha,\omega,\omega_1,\omega_2$ 是 $\mathbb{R}^n$ 上的形式, 则
\begin{enumerate}[label=(\roman*),leftmargin=*]
\item $\mathrm{d}(\omega_1+\omega_2)=\mathrm{d}\omega_1+\mathrm{d}\omega_2$;
\item 若 $\omega$ 是 $k$-形式, 则 $\mathrm{d}(\omega\wedge\alpha)=\mathrm{d}\omega\wedge\alpha+(-1)^k\omega\wedge\mathrm{d}\alpha$;
\item $\mathrm{d}(\mathrm{d}\omega)=0$;
\item 若 $f:\mathbb{R}^m\to\mathbb{R}^n$ 光滑, 则 $\mathrm{d}(f^*\omega)=f^*(\mathrm{d}\omega)$.
\end{enumerate}
\end{RiemannProposition}

\begin{RiemannProof}
性质 (i) 显然成立. 利用 (i), 只需对 $\omega=a_I\,\mathrm{d}x^I$ 和 $\alpha=b_J\,\mathrm{d}x^J$ 证明 (ii):
\begin{align*}
  \mathrm{d}(\omega\wedge\alpha)
  &=\mathrm{d}(a_Ib_J\,\mathrm{d}x^I\wedge\mathrm{d}x^J)
    =\mathrm{d}(a_Ib_J)\wedge\mathrm{d}x^I\wedge\mathrm{d}x^J\\
  &=(b_J\,\mathrm{d}a_I+a_I\,\mathrm{d}b_J)\wedge\mathrm{d}x^I\wedge\mathrm{d}x^J\\
  &=b_J\,\mathrm{d}a_I\wedge\mathrm{d}x^I\wedge\mathrm{d}x^J
    +a_I\,\mathrm{d}b_J\wedge\mathrm{d}x^I\wedge\mathrm{d}x^J\\
  &=\mathrm{d}\omega\wedge\alpha+(-1)^ka_I\,\mathrm{d}x^I\wedge\mathrm{d}b_J\wedge\mathrm{d}x^J\\
  &=\mathrm{d}\omega\wedge\alpha+(-1)^k\omega\wedge\mathrm{d}\alpha.
\end{align*}
同样, 为证明 (iii), 只需考虑形式 $\omega=a_I\,\mathrm{d}x^I$. 由于
\[
  \mathrm{d}\omega=\mathrm{d}a_I\wedge\mathrm{d}x^I
  =\sum_{i=1}^n\frac{\partial a_I}{\partial x^i}\,\mathrm{d}x^i\wedge\mathrm{d}x^I,
\]
有
\begin{align*}
  \mathrm{d}(\mathrm{d}\omega)
  &=\sum_{j=1}^n\sum_{i=1}^n\frac{\partial^2a_I}{\partial x^j\partial x^i}
    \,\mathrm{d}x^j\wedge\mathrm{d}x^i\wedge\mathrm{d}x^I\\
  &=\sum_{i=1}^n\sum_{j<i}\left(\frac{\partial^2a_I}{\partial x^j\partial x^i}
    -\frac{\partial^2a_I}{\partial x^i\partial x^j}\right)
    \mathrm{d}x^j\wedge\mathrm{d}x^i\wedge\mathrm{d}x^I=0.
\end{align*}

为证明 (iv), 先考虑 0-形式 $g$:
\begin{align*}
  f^*(\mathrm{d}g)
  &=f^*\left(\sum_{i=1}^n\frac{\partial g}{\partial x^i}\,\mathrm{d}x^i\right)
    =\sum_{i=1}^n\left(\frac{\partial g}{\partial x^i}\circ f\right)\mathrm{d}f^i\\
  &=\sum_{i,j=1}^n\left(\left(\frac{\partial g}{\partial x^i}\circ f\right)
    \frac{\partial f^i}{\partial x^j}\right)\mathrm{d}x^j\\
  &=\sum_{j=1}^n\frac{\partial(g\circ f)}{\partial x^j}\,\mathrm{d}x^j
    =\mathrm{d}(g\circ f)=\mathrm{d}(f^*g).
\end{align*}
再取 $\omega=a_I\,\mathrm{d}x^I$, 有
\begin{align*}
  \mathrm{d}(f^*\omega)&=\mathrm{d}((f^*a_I)\,\mathrm{d}f^I)\\
  &=\mathrm{d}(f^*a_I)\wedge\mathrm{d}f^I+(f^*a_I)\,\mathrm{d}(\mathrm{d}f^I)
    =\mathrm{d}(f^*a_I)\wedge\mathrm{d}f^I\\
  &=(f^*\mathrm{d}a_I)\wedge(f^*\mathrm{d}x^I)
    =f^*(\mathrm{d}a_I\wedge\mathrm{d}x^I)=f^*(\mathrm{d}\omega)
\end{align*}
(其中 $\mathrm{d}f^I$ 表示形式 $\mathrm{d}f^{i_1}\wedge\cdots\wedge\mathrm{d}f^{i_k}$), 由此得证.
\end{RiemannProof}

现在设 $\omega$ 是光滑流形 $M$ 上的微分 $k$-形式. 定义 $(k+1)$-形式 $\mathrm{d}\omega$ 为这样一个光滑形式: 对每个参数化 $\varphi_\alpha:U_\alpha\to M$, 它的局部表示为 $\mathrm{d}\omega_\alpha$, 其中 $\omega_\alpha:=\varphi_\alpha^*\omega$ 是 $\omega$ 的局部表示, 即在 $\varphi_\alpha(U)$ 上有 $\mathrm{d}\omega=(\varphi_\alpha^{-1})^*(\mathrm{d}\omega_\alpha)$. 给定另一个参数化 $\varphi_\beta:U_\beta\to M$, 使得 $W:=\varphi_\alpha(U_\alpha)\cap\varphi_\beta(U_\beta)\ne\varnothing$, 则容易验证
\[
  (\varphi_\alpha^{-1}\circ\varphi_\beta)^*\omega_\alpha=\omega_\beta.
\]
令 $f$ 等于 $\varphi_\alpha^{-1}\circ\varphi_\beta$, 有
\[
  f^*(\mathrm{d}\omega_\alpha)=\mathrm{d}(f^*\omega_\alpha)=\mathrm{d}\omega_\beta.
\]
于是,
\begin{align*}
  (\varphi_\beta^{-1})^*\mathrm{d}\omega_\beta
  &=(\varphi_\beta^{-1})^*f^*(\mathrm{d}\omega_\alpha)\\
  &=(f\circ\varphi_\beta^{-1})^*(\mathrm{d}\omega_\alpha)\\
  &=(\varphi_\alpha^{-1})^*(\mathrm{d}\omega_\alpha),
\end{align*}
所以两个定义在重叠集 $W$ 上一致. 因此, $\mathrm{d}\omega$ 是良定义的. 光滑流形上形式的外微分也满足 Proposition~\ref{prop:properties-of-exterior-derivative} 的各条性质, 这一点留作练习.

\begin{RiemannExercise}\label{ex:differential-forms}
\begin{enumerate}[label=(\arabic*),ref=\arabic*,leftmargin=*]
\item\label{ex:prove-pullback-of-differential-forms} 证明 Proposition~\ref{prop:pullback-of-differential-forms}.
\item\label{ex:intrinsic-exterior-derivative} (\emph{外微分}) 设 $M$ 是光滑流形. 给定 $M$ 上的 $k$-形式 $\omega$, 可以不使用局部坐标定义其外微分 $\mathrm{d}\omega$: 给定 $k+1$ 个向量场 $X_1,\ldots,X_{k+1}\in\mathfrak{X}(M)$, 定义
\begin{align*}
  \mathrm{d}\omega(X_1,\ldots,X_{k+1})
  &:=\sum_{i=1}^{k+1}(-1)^{i-1}X_i\cdot\omega(X_1,\ldots,\widehat X_i,\ldots,X_{k+1})\\
  &\quad+\sum_{i<j}(-1)^{i+j}\omega([X_i,X_j],X_1,\ldots,\widehat X_i,\ldots,\widehat X_j,\ldots,X_{k+1}),
\end{align*}
其中帽号表示省略该变量.
\begin{enumerate}[label=(\alph*),leftmargin=*]
\item 证明上述定义的 $\mathrm{d}\omega$ 确实是 $M$ 上的 $(k+1)$-形式, 即
\begin{enumerate}[label=(\roman*),leftmargin=*]
\item $\mathrm{d}\omega(X_1,\ldots,X_i+Y_i,\ldots,X_{k+1})={}$\\
$\mathrm{d}\omega(X_1,\ldots,X_i,\ldots,X_{k+1})+\mathrm{d}\omega(X_1,\ldots,Y_i,\ldots,X_{k+1})$;
\item 对任意可微函数 $f$, 有 $\mathrm{d}\omega(X_1,\ldots,fX_j,\ldots,X_{k+1})=f\,\mathrm{d}\omega(X_1,\ldots,X_{k+1})$;
\item $\mathrm{d}\omega$ 是交错的;
\item $\mathrm{d}\omega(X_1,\ldots,X_{k+1})(p)$ 只依赖于 $(X_1)_p,\ldots,(X_{k+1})_p$.
\end{enumerate}
\item 设 $x:W\to\mathbb{R}^n$ 是 $M$ 的坐标系, 且 $\omega=\sum_Ia_I\,\mathrm{d}x^{i_1}\wedge\cdots\wedge\mathrm{d}x^{i_k}$ 是 $\omega$ 在这些坐标下的表达式 (其中 $a_I$ 为光滑函数). 证明 $\mathrm{d}\omega$ 的局部表达式与外微分的局部定义所用的表达式相同, 即
\[
  \mathrm{d}\omega=\sum_I\mathrm{d}a_I\wedge\mathrm{d}x^{i_1}\wedge\cdots\wedge\mathrm{d}x^{i_k}.
\]
\end{enumerate}
\item\label{ex:properties-of-exterior-derivative-on-manifolds} 证明对光滑流形上的形式定义的外微分满足 Proposition~\ref{prop:properties-of-exterior-derivative} 的性质.
\item\label{ex:exterior-derivative-curl-and-divergence} 证明:
\begin{enumerate}[label=(\alph*),leftmargin=*]
\item 若 $\omega=f^1\,\mathrm{d}x+f^2\,\mathrm{d}y+f^3\,\mathrm{d}z$ 是 $\mathbb{R}^3$ 上的 1-形式, 则
\[
  \mathrm{d}\omega=g^1\,\mathrm{d}y\wedge\mathrm{d}z
    +g^2\,\mathrm{d}z\wedge\mathrm{d}x+g^3\,\mathrm{d}x\wedge\mathrm{d}y,
\]
其中 $(g^1,g^2,g^3)=\operatorname{curl}(f^1,f^2,f^3)$;
\item 若 $\omega=f^1\,\mathrm{d}y\wedge\mathrm{d}z+f^2\,\mathrm{d}z\wedge\mathrm{d}x+f^3\,\mathrm{d}x\wedge\mathrm{d}y$ 是 $\mathbb{R}^3$ 上的 2-形式, 则
\[
  \mathrm{d}\omega=\operatorname{div}(f^1,f^2,f^3)\,\mathrm{d}x\wedge\mathrm{d}y\wedge\mathrm{d}z.
\]
\end{enumerate}
\item\label{ex:de-rham-cohomology} (\emph{De Rham 上同调}) 若 $\mathrm{d}\omega=0$, 则称 $k$-形式 $\omega$ \textbf{闭}. 若存在 $(k-1)$-形式 $\beta$, 使得 $\omega=\mathrm{d}\beta$, 则称 $\omega$ \textbf{恰当}. 注意, 每个恰当形式都是闭形式. 设 $Z^k$ 是 $M$ 上所有闭 $k$-形式的集合, 在 $Z^k$ 上的形式之间定义如下关系: $\alpha\sim\beta$ 当且仅当它们相差一个恰当形式, 即对某个 $(k-1)$-形式 $\theta$, 有 $\beta-\alpha=\mathrm{d}\theta$.
\begin{enumerate}[label=(\alph*),leftmargin=*]
\item 证明该关系是等价关系.
\item 设 $H^k(M)$ 是相应的等价类集合 (称为 $M$ 的 \textbf{$k$ 维 de Rham 上同调空间}). 证明形式的加法和数乘确实在 $H^k(M)$ 上定义了向量空间结构.
\item 设 $f:M\to N$ 是光滑映射. 证明:
\begin{enumerate}[label=(\roman*),leftmargin=*]
\item 拉回 $f^*$ 将闭形式映为闭形式, 将恰当形式映为恰当形式;
\item 若 $\alpha\sim\beta$ 在 $N$ 上成立, 则 $f^*\alpha\sim f^*\beta$ 在 $M$ 上成立;
\item $f^*$ 诱导上同调上的线性映射 $f^\sharp:H^k(N)\to H^k(M)$, 由 $f^\sharp[\omega]=[f^*\omega]$ 自然定义;
\item 若 $g:L\to M$ 是另一个光滑映射, 则 $(f\circ g)^\sharp=g^\sharp\circ f^\sharp$.
\end{enumerate}
\item 证明 $H^0(M)$ 的维数等于 $M$ 的连通分支数.
\item 证明对每个 $k>\dim M$, 有 $H^k(M)=0$.
\end{enumerate}
\item\label{ex:homotopy-operator-and-poincare-lemma} 设 $M$ 是 $n$ 维流形, $U$ 是 $\mathbb{R}^n$ 的开子集, $\omega$ 是 $\mathbb{R}\times U$ 上的 $k$-形式. 将 $\omega$ 写成
\[
  \omega=\mathrm{d}t\wedge\sum_Ia_I\,\mathrm{d}x^I+\sum_Jb_J\,\mathrm{d}x^J,
\]
其中 $I=(i_1,\ldots,i_{k-1})$ 和 $J=(j_1,\ldots,j_k)$ 为递增指标序列, $(x^1,\ldots,x^n)$ 为 $U$ 中的坐标, $t$ 为 $\mathbb{R}$ 中的坐标. 考虑由下式定义的算子 $Q$:
\[
  Q(\omega)_{(t,x)}=\sum_I\left(\int_{t_0}^ta_I\,\mathrm{d}s\right)\mathrm{d}x^I,
\]
它将 $\mathbb{R}\times U$ 中的 $k$-形式 $\omega$ 变为 $(k-1)$-形式.
\begin{enumerate}[label=(\alph*),ref=\alph*,leftmargin=*]
\item\label{ex:homotopy-operator-coordinate-naturality} 设 $f:V\to U$ 是 $\mathbb{R}^n$ 中开子集之间的微分同胚. 证明诱导的微分同胚 $\widetilde f:=\operatorname{id}\times f:\mathbb{R}\times V\to\mathbb{R}\times U$ 满足
\[
  \widetilde f^*\circ Q=Q\circ\widetilde f^*.
\]
\item\label{ex:homotopy-operator-on-manifolds} 利用 (\ref{ex:homotopy-operator-coordinate-naturality}), 构造一个算子 $Q$, 将 $\mathbb{R}\times M$ 上的 $k$-形式映为 $(k-1)$-形式, 并使得对任意微分同胚 $f:M\to N$, 诱导的微分同胚 $\widetilde f:=\operatorname{id}\times f:\mathbb{R}\times M\to\mathbb{R}\times N$ 满足 $\widetilde f^*\circ Q=Q\circ\widetilde f^*$. 证明此算子是线性的.
\item 考虑 (\ref{ex:homotopy-operator-on-manifolds}) 中定义的算子 $Q$, 以及 $M$ 在\textquotedblleft 层\textquotedblright $t_0$ 处的包含映射 $i_{t_0}:M\to\mathbb{R}\times M$, 它由 $i_{t_0}(p)=(t_0,p)$ 定义. 证明 $\omega-\pi^*i_{t_0}^*\omega=\mathrm{d}Q\omega+Q\mathrm{d}\omega$, 其中 $\pi:\mathbb{R}\times M\to M$ 是到 $M$ 的投影.
\item\label{ex:homotopy-operator-cylinder-cohomology} 证明映射 $\pi^\sharp:H^k(M)\to H^k(\mathbb{R}\times M)$ 和 $i_{t_0}^\sharp:H^k(\mathbb{R}\times M)\to H(M)$ 互为逆映射 (因而 $H^k(M)$ 同构于 $H^k(\mathbb{R}\times M)$).
\item 利用 (\ref{ex:homotopy-operator-cylinder-cohomology}) 证明: 对 $k>0$ 和 $n>0$, $\mathbb{R}^n$ 中每个闭 $k$-形式都是恰当的, 即当 $k>0$ 时 $H^k(\mathbb{R}^n)=0$.
\item 利用 (\ref{ex:homotopy-operator-cylinder-cohomology}) 证明: 若 $f,g:M\to N$ 是光滑流形之间的两个\textbf{光滑同伦的映射} (意即存在光滑映射 $H:\mathbb{R}\times M\to N$, 使得对某些固定的 $t_0,t_1\in\mathbb{R}$, 有 $H(t_0,p)=f(p)$ 和 $H(t_1,p)=g(p)$), 则 $f^\sharp=g^\sharp$.
\item 若恒等映射 $\operatorname{id}:M\to M$ 与常值映射光滑同伦, 则称 $M$ \textbf{可缩}. 证明 $\mathbb{R}^n$ 可缩.
\item (\emph{Poincaré 引理}) 设 $M$ 是可缩的光滑流形. 证明 $M$ 上每个闭形式都是恰当的, 即对所有 $k>0$, 有 $H^k(M)=0$.

{\small (\textbf{Remark:} 本题基于 \RGcite{ch02-GP73} 中的一道练习).}
\end{enumerate}
\item\label{ex:lie-derivative-of-differential-form} (\emph{微分形式的 Lie 导数}) 给定向量场 $X\in\mathfrak{X}(M)$, 定义\textbf{形式 $\omega$ 沿 $X$ 的 Lie 导数}为
\[
  L_X\omega:=\left.\frac{\mathrm{d}}{\mathrm{d}t}(\psi_t^*\omega)\right|_{t=0},
\]
其中 $\psi_t=F(\cdot,t)$, $F$ 是 $X$ 在 $p$ 处的局部流 [参见 Exercise~\ref{ex:tensor-fields}(\ref{ex:lie-derivative-of-tensors})]. 证明 Lie 导数满足下列性质:
\begin{enumerate}[label=(\alph*),leftmargin=*]
\item $L_X(\omega_1\wedge\omega_2)=(L_X\omega_1)\wedge\omega_2+\omega_1\wedge(L_X\omega_2)$;
\item $\mathrm{d}(L_X\omega)=L_X(\mathrm{d}\omega)$;
\item \textbf{Cartan 公式:} $L_X\omega=\iota(X)\mathrm{d}\omega+\mathrm{d}(\iota(X)\omega)$;
\item $L_X(\iota(Y)\omega)=\iota(L_XY)\omega+\iota(Y)L_X\omega$\\\relax
[参见第~\ref{chap:differentiable-manifolds}~章的 Exercise~\ref{ex:vector-fields}(\ref{ex:lie-derivative-of-vector-field}) 和 Exercise~\ref{ex:tensors}(\ref{ex:contraction-of-alternating-tensors})].
\end{enumerate}
\end{enumerate}
\end{RiemannExercise}
